\section{Badania}

Jedna z metod badań społecznych, w której do zbierania informacji od respondentów wykorzystuje się wystandaryzowany kwestionariusz. \cite{French2010a}

Dla ustalenia badanej populacji konsumentów \cite{nbp2012} stosuje się dobór losowy, gdy niewielka jest ich liczba i nieznana ich struktura, dobór losowy systematyczny \cite{nbp2012}, gdy dostępna jest lista całej populacji oraz dobór losowy \cite{systems2001} warstwowy dla uzyskania bardziej dokładnych wyników charakteryzujących badanych. \cite{Shotton2000, French2010, Shotton2010, Shotton1989, French2010a, Shotton1999}

Stan negatywnych emocji i frustracji w niektórych kręgach akademickiej makroekonomii \textcite[25]{krugman2009} najlepiej oddaje skrajna wypowiedź Krugmana \parencite*{krugman2009}: „Większość z tego, co stworzono w ramach makroekonomii w ciągu ostatnich 30 lat okazało się w najlepszym razie bezużyteczne lub w najgorszym razie szkodliwe” \citepage[20]{krugman2009}.

\begin{table}[h!]
	\begin{center}
		\caption{{\bfseries Wyniki badań}}
		\label{tab:table1}
		\begin{tabular}{l|c|r} % <-- Alignments: 1st column left, 2nd middle and 3rd right, with vertical lines in between
			\textbf{Value 1} & \textbf{Value 2} & \textbf{Value 3}\\
			$\alpha$ & $\beta$ & $\gamma$ \\
			\hline
			1 & 1110.1 & a\\
			2 & 10.1 & b\\
			3 & 23.113231 & c\\
		\end{tabular}
	\end{center}
	\begin{center}
	    \footnotesize{Źródło: \protect\cite{Shotton1999}}
	\end{center}
\end{table}